% !TEX root = ../ApproxDA-TransUNet.tex
\begin{abstract}
Dual-attention mechanisms have demonstrated strong performance in medical image segmentation by jointly modeling long-range spatial and channel dependencies, but their quadratic complexity has motivated numerous approximation methods whose task-dependent behavior remains poorly understood. To bridge this gap, we propose Approximate Dual Attention TransUNet (ApproxDA-TransUNet), an efficient dual-attention approximation framework for hybrid CNN-Transformer architectures that replaces global attention with controllable windowed low-rank spatial attention and grouped channel attention. Unlike existing efficient attention methods that primarily pursue computational savings, ApproxDA-TransUNet also provides a controllable framework for systematically studying how attention approximation affects segmentation performance. Extensive experiments across five representative benchmarks (Synapse, ACDC, Kvasir-SEG, CVC-ClinicDB, and ISIC 2018) demonstrate the best performance among compared methods with consistent improvements over DA-TransUNet: 80.94\%, 88.97\%, 90.17\%, 90.99\%, and 89.58\% Dice on Synapse, ACDC, Kvasir-SEG, CVC-ClinicDB, and ISIC 2018, respectively. These results also reveal two key observations: First, we identify \textbf{Global Context Sensitivity (GCS)}: tasks requiring long-range dependencies are far more sensitive to approximation scale than texture-dominated ones, and inter-class spatial separation is the strongest observed correlate of this sensitivity among the factors we examined. Second, we observe \textbf{Gate Collapse}: in our evaluated configurations, the learnable attention gate settles into near-uniform routing, which a gradient-based fixed-point analysis suggests is associated with gradient symmetry between branches, undermining adaptive fusion. These findings reframe approximation as a task-dependent inductive bias: rather than a pure efficiency tool, it can improve segmentation accuracy when the structural prior induced by spatial restriction, low-rank projection, or both aligns with the task's contextual requirements. Our code and analysis offer both a strong segmentation baseline and actionable principles for designing efficient medical transformers.
\end{abstract}

\begin{IEEEkeywords}
medical image segmentation, dual attention, efficient transformers, attention approximation, global context sensitivity, gate collapse
\end{IEEEkeywords}
